% Dòng thời gian biểu diễn tri thức và RAG: mạng ngữ nghĩa -> Web ngữ nghĩa -> Google KG -> Transformer -> RAG -> ReAct -> GraphRAG.
\begin{tikzpicture}[font=\small, x=1cm, y=1cm,
  ev/.style={hop=#1, text width=2.85cm, minimum height=1.55cm, font=\scriptsize},
  dot/.style={circle, fill=#1, draw=white, line width=1.5pt, minimum size=4.5mm, inner sep=0pt}]
  \shade[left color=cSky!60, right color=cPurple!70, rounded corners=4pt] (0,-0.22) rectangle (16.6,0.22);
  \draw[-{Stealth[length=4mm]}, line width=3pt, cPurple!80!black] (16.4,0) -- (17.2,0);
  \foreach \x/\yr/\c/\t/\up in {
      0.8/1968/cSlate/{\textbf{Mạng ngữ nghĩa}\\Quillian: khái niệm là nút, quan hệ là cạnh}/1,
      3.2/1993/cBrown/{\textbf{Ontology}\\Gruber: đặc tả tường minh một khái niệm hoá}/0,
      5.6/2001/cGreen/{\textbf{Web ngữ nghĩa}\\RDF, RDFS, OWL, SPARQL (W3C)}/1,
      8.0/2012/cBlue/{\textbf{Google Knowledge Graph}\\``things, not strings''}/0,
      10.4/2017/cOrange/{\textbf{Transformer}\\nền tảng của LLM và embedding}/1,
      12.8/2020/cRed/{\textbf{RAG}\\Lewis et al.: bộ nhớ tham số + phi tham số}/0,
      15.2/2023/cPurple/{\textbf{ReAct, GraphRAG}\\Yao 2023, Edge 2024: suy luận + hành động, cộng đồng}/1} {
    \node[dot=\c] at (\x,0) {};
    \node[font=\small\bfseries, text=\c!75!black] at (\x, {0.62-1.24*\up}) {\yr};
    \ifnum\up=1
      \node[ev=\c, anchor=south] (b) at (\x,1.0) {\t};
      \draw[thick, \c] (\x,0.25) -- (b.south);
    \else
      \node[ev=\c, anchor=north] (b) at (\x,-1.0) {\t};
      \draw[thick, \c] (\x,-0.25) -- (b.north);
    \fi
  }
\end{tikzpicture}
